\section{Related Work}
\label{sec:related}

\subsection{Knowledge Tracing Models and Benchmarking}
Knowledge Tracing (KT) formulates the task of modeling a learner's mastery over time to predict their probability of correctly answering future questions based on their interaction history \cite{abdelrahman2023knowledge}. Classical approaches, such as Item Response Theory (IRT) \cite{rasch1960probabilistic,embretson2000item} and Bayesian Knowledge Tracing (BKT) \cite{corbett1994knowledge}, model static ability and latent mastery states, respectively. The introduction of Deep Knowledge Tracing (DKT) \cite{piech2015deep} shifted the paradigm toward recurrent neural networks, later followed by memory-augmented structures like DKVMN \cite{zhang2017dynamic} and attention-based models such as AKT \cite{ghosh2020context} and simpleKT \cite{liu2023simplekt}. Recently, benchmarking toolkits such as pyKT \cite{liu2022pykt} have significantly improved the standardization of data preprocessing and model evaluation. In this paper, we focus on diagnostics rather than proposing a new KT architecture; we utilize representative baselines to validate our protocols.

\subsection{Graph-Augmented Knowledge Tracing}
Recent advances in KT often incorporate auxiliary information using graph neural networks to model prerequisite, similarity, co-occurrence, or concept dependency relationships between knowledge components (KCs) \cite{abdelrahman2023knowledge}. Models like GKT \cite{nakagawa2019graph} and GIKT \cite{yang2021gikt} construct KC-item or KC-KC dependency graphs, while contrastive/self-supervised approaches like CL4KT \cite{lee2022cl4kt} leverage structurally augmented views. While these graph-augmented architectures effectively capture rich relations, they introduce significant complexity. Crucially, graph and self-supervised KT models require highly reliable evaluation protocols to prevent leakage and to report performance across varying KC-frequency strata. To maintain a focused evaluation on fundamental sequence modeling and calibration, our diagnostic protocol evaluates standard sequential KT models, establishing a rigorous foundation that can be extended to graph-based architectures in the future.

\subsection{Sparse and Cold-start Knowledge Components}
Educational datasets naturally exhibit long-tail distributions where a small fraction of dense KCs dominate interactions, leaving a vast majority of KCs with sparse training evidence \cite{choi2020ednet,liu2023xes3g5m}. It is important to distinguish sparse and cold-start conditions in KT from typical cold-start scenarios in recommender systems. While recommender systems typically define cold-start at the user or item level (i.e., new users or items), KT sparse/cold-start issues relate to concepts (KCs) that have insufficient interaction history in the training data \cite{pelanek2015modeling,choffin2019das3h}. Our work utilizes train-only KC-frequency stratification to group KCs into very sparse, sparse, medium, and dense buckets. This systematic stratification is designed to ensure that overall AUC metrics do not obscure model degradation on data-poor concepts.

\subsection{Calibration of Probabilistic Predictions}
Model calibration is an essential property for KT models because their outputs---the predicted probability of correctness---can directly drive high-stakes educational interventions. While deep neural networks can achieve high classification accuracy (e.g., high AUC), they may remain poorly calibrated, often producing overconfident predictions \cite{guo2017calibration}. In machine learning literature, calibration error is commonly quantified via Expected Calibration Error (ECE) \cite{naeini2015obtaining}. Furthermore, the Brier score \cite{brier1950verification} and its decomposition into uncertainty, reliability, and resolution \cite{degroot1983comparison} separate the inherent task difficulty from the model's predictive calibration. In this paper, we report ECE, the Brier score, its decomposition, and reliability diagrams. Crucially, we evaluate calibration not only overall but also across KC-frequency strata to pinpoint reliability failures under sparsity.

\subsection{Reproducibility, Leakage Control, and Statistical Testing}
The rapid proliferation of ML models has raised widespread concerns regarding reproducibility and data leakage \cite{kapoor2023leakage}, particularly in sequence modeling tasks like KT \cite{gervet2020deep}. Strict leakage control is necessary to prevent artificial inflation of performance. We identify and mitigate several leakage vectors, including split leakage, preprocessing leakage, KC mapping leakage, sparse-bucket leakage, calibration leakage, hyperparameter leakage, and cold-start leakage. For comparing competitive models, DeLong's test provides a standard approach for correlated ROC curves \cite{delong1988comparing}, and the Wilcoxon signed-rank test is appropriate for comparing repeated measurements \cite{wilcoxon1945individual}. Statistical testing is used only where prediction-level outputs permit paired comparison; otherwise, we report descriptive diagnostics with sample-size-aware interpretation.
